\section{From Description to Intervention}

The previous chapters argued that behavioural failure is often a state-transition
problem rather than a simple shortage of desire. That claim matters only if it
changes design practice. An intervention chapter therefore has a stricter job
than a motivational chapter: it must specify which control variable is being
changed, why that variable is cheaper or more reliable than direct exhortation,
and which parts of the design are supported by stronger evidence rather than
mere intuition.

In this book, intervention design is governed by five ideas:
\begin{enumerate}[nosep]
  \item interventions act on transitions, not on abstract identity claims;
  \item decision windows matter because leverage is timing-sensitive;
  \item repeated exposure and history dependence matter because frameworks are
        not rebuilt from zero at every attempt;
  \item attractor and metastability language helps explain why some local
        redesigns have outsized downstream effects;
  \item source discipline matters because empirical findings, theoretical
        mechanism-building, and speculative frontier claims should not be mixed
        as if they had equal status.
\end{enumerate}
